\documentclass[11pt]{article}
\usepackage{amsmath, amssymb, amsthm}
\usepackage{geometry}
\usepackage{hyperref}
\usepackage{array}
\geometry{margin=1in}

\newtheorem{theorem}{Theorem}
\newtheorem{lemma}[theorem]{Lemma}
\newtheorem{proposition}[theorem]{Proposition}
\newtheorem{corollary}[theorem]{Corollary}
\newtheorem{conjecture}[theorem]{Conjecture}
\theoremstyle{definition}
\newtheorem{definition}[theorem]{Definition}
\newtheorem{observation}[theorem]{Observation}
\newtheorem{remark}[theorem]{Remark}

\title{The Eigenvalue Rosetta Stone at $n=4$:\\
Explicit Verification}
\author{Clio}
\date{2026-04-22}

\begin{document}
\maketitle

\begin{abstract}
We verify the conjectural three-factor decomposition
$e_\lambda(q) = q^{n(\lambda)} \cdot (1+q)^{\lfloor n/2\rfloor}
\cdot p_\lambda(q)$
for the per-irrep eigenvalues of the staircase Hecke product
$\Pi_q = (1+T_1)(1+T_2)(1+T_1)(1+T_3)(1+T_2)(1+T_1) \in H_q(S_4)$ acting
on each Specht module $V_\lambda$ with $\lambda \vdash 4$. For each of the
three irreps $\lambda \in \{(4),(3,1),(2,2)\}$ where $\Pi_q$ is nonzero, the
residual polynomial $p_\lambda(q) \in \mathbb{Z}_{\geq 0}[q]$ is palindromic,
with coefficients $(1,4,6,4,1)$, $(1,4,1)$, and $(1)$ respectively. The
three-factor form instantiates the ``Eigenvalue Rosetta Stone'' hypothesis
on $n=4$: the statistic $n(\lambda)$ (the cocharge-like piece), the uniform
factor $(1+q)^{\lfloor n/2\rfloor}$ (the descent-paired piece), and the
palindromic remainder $p_\lambda$ (the Kazhdan--Lusztig piece) each arise
from a distinct structural source. On $V_{(2,1,1)}$ and $V_{(1^4)}$ the
product $\Pi_q$ acts as zero, consistent with the rank-hierarchy prediction
that $V_\lambda$ survives under $\Pi_q$ iff $\ell(\lambda) \leq (n+1)/2$.
We close with a conjecture tying the six nonzero eigenvalues (with
multiplicity) to the six indecomposable Soergel summands of the
Bott--Samelson object $BS(\mathbf{i})$.
\end{abstract}

\section{Setup}

\begin{definition}[Staircase product]\label{def:pi}
Let $H_q(S_4)$ be the Iwahori--Hecke algebra with quadratic relation
$T_i^2 = (q-1) T_i + q$. The \emph{staircase Hecke product} is
\[
\Pi_q \;:=\; F_1 F_2 F_1 F_3 F_2 F_1,
\qquad F_i := 1 + T_i,
\]
corresponding to the reduced expression
$\mathbf{i} = s_1 s_2 s_1 s_3 s_2 s_1$ for $w_0 \in S_4$. For each
partition $\lambda \vdash 4$, let $V_\lambda$ denote the Specht
$H_q(S_4)$-module (i.e., the deformation of the $S_4$-Specht module) and
$\rho_\lambda \colon H_q(S_4) \to \mathrm{End}(V_\lambda)$ its representation.
\end{definition}

\begin{definition}[Statistic $n(\lambda)$]\label{def:nlambda}
For $\lambda = (\lambda_1, \lambda_2, \ldots) \vdash n$, define
$n(\lambda) := \sum_{i\geq 1} (i-1)\,\lambda_i$. Equivalently,
$n(\lambda) = \sum_{\square\in\lambda}\mathrm{leg}(\square) =
\sum_{i} \binom{\lambda'_i}{2}$ where $\lambda'$ is the conjugate.
\end{definition}

\begin{definition}[Descent-paired piece]
The descent-paired subset
$D_n := \{ w \in S_n : w(2i-1) > w(2i) \; \forall\, 1 \leq i \leq \lfloor n/2\rfloor\}$
has $|D_n| = n!/2^{\lfloor n/2\rfloor}$. For $n=4$, $|D_4| = 24/4 = 6$.
\end{definition}

The central object of this note is the polynomial identity
\begin{equation}\label{eq:rosetta}
e_\lambda(q) \;=\; q^{n(\lambda)} \cdot (1+q)^{\lfloor n/2 \rfloor}
\cdot p_\lambda(q)
\quad\text{for each}\; \lambda\vdash n=4,
\end{equation}
where $e_\lambda(q)$ is the nonzero eigenvalue of $\rho_\lambda(\Pi_q)$ on
a copy of $V_\lambda$ (when $\rho_\lambda(\Pi_q) \ne 0$), and $p_\lambda$
is a residual polynomial.

\begin{theorem}[Rosetta decomposition at $n=4$]\label{thm:main}
For each $\lambda \vdash 4$, the following table gives the dimension of
$V_\lambda$, the statistic $n(\lambda)$, the rank of $\rho_\lambda(\Pi_q)$,
the eigenvalue $e_\lambda(q)$ of $\rho_\lambda(\Pi_q)$ on a copy of
$V_\lambda$ (when nonzero), and the residual polynomial
$p_\lambda(q) := e_\lambda(q) / \bigl(q^{n(\lambda)}(1+q)^{2}\bigr)$:
\begin{center}
\renewcommand{\arraystretch}{1.25}
\begin{tabular}{c|c|c|c|l|l}
$\lambda$ & $\dim V_\lambda$ & $n(\lambda)$ & rank & $e_\lambda(q)$ & $p_\lambda(q)$ \\\hline
$(4)$     & $1$ & $0$ & $1$ & $(1+q)^6$             & $(1+q)^4 = 1+4q+6q^2+4q^3+q^4$ \\
$(3,1)$   & $3$ & $1$ & $1$ & $q(1+q)^2(q^2+4q+1)$  & $q^2+4q+1$ \\
$(2,2)$   & $2$ & $2$ & $1$ & $q^2(1+q)^2$          & $1$ \\
$(2,1,1)$ & $3$ & $3$ & $0$ & $0$                   & --- (rank 0) \\
$(1,1,1,1)$ & $1$ & $6$ & $0$ & $0$                 & --- (rank 0) \\
\end{tabular}
\end{center}
In particular, identity~\eqref{eq:rosetta} holds for all $\lambda \vdash 4$
with $\rho_\lambda(\Pi_q) \ne 0$, and each $p_\lambda(q)$ is a palindromic
polynomial with nonnegative integer coefficients.
\end{theorem}

\section{Proof of Theorem~\ref{thm:main}}

We compute the characteristic polynomial of left multiplication by
$\Pi_q$ on the regular representation, factor it by irrep, and check
identity~\eqref{eq:rosetta} term by term.

\begin{lemma}[Characteristic polynomial on the regular representation]
\label{lem:charpoly}
Let $M_{\Pi_q} \in \mathrm{Mat}_{24}(\mathbb{Z}[q])$ be the matrix of left
multiplication by $\Pi_q$ on the length-lex $T$-basis of $H_q(S_4)$. Then
\begin{align*}
\chi_{M_{\Pi_q}}(x) \;=\; x^{18} &\cdot \bigl(x - (q^4+2q^3+q^2)\bigr)^{2}
\cdot \bigl(x - (q^5+6q^4+10q^3+6q^2+q)\bigr)^{3}\\
&\cdot\bigl(x - (q^6+6q^5+15q^4+20q^3+15q^2+6q+1)\bigr).
\end{align*}
The three nonzero factors correspond to $V_{(2,2)}$, $V_{(3,1)}$, $V_{(4)}$
respectively; the $x^{18}$ factor absorbs $V_{(2,1,1)}$ (dim $3$, mult $3$)
and $V_{(1^4)}$ (dim $1$, mult $1$), plus the kernels inside each nonzero
block.
\end{lemma}

\begin{proof}
Direct symbolic computation via SymPy's \texttt{.charpoly().factor()}; see
\texttt{/home/clio/projects/bms-test/followup\_structure.py}
(\texttt{charpoly\_factored} field of \texttt{structure\_result.pickle}).
Assignment to irreps via isotypic-projection rank computation in
\texttt{compare\_pi\_bms.py} (\texttt{per\_irrep\_ranks\_q1} field of
\texttt{pi\_bms\_result.pickle}): the ranks of $\rho_\lambda(\Pi_q)$ on a
single copy of $V_\lambda$ at $q=1$ are
$\{(4)\mapsto 1, (3,1)\mapsto 1, (2,2)\mapsto 1, (2,1,1)\mapsto 0, (1^4)\mapsto 0\}$.
Generic rank equals the rank at $q=1$ by lower-semicontinuity plus attainment
of the upper bound at $q \in \{2,3,5\}$ (separately verified).
\end{proof}

\begin{lemma}[Factorization of the nonzero blocks]\label{lem:factor}
In $\mathbb{Z}[q]$,
\begin{align*}
q^6+6q^5+15q^4+20q^3+15q^2+6q+1 &= (1+q)^6,\\
q^5+6q^4+10q^3+6q^2+q &= q(1+q)^2 (q^2+4q+1),\\
q^4+2q^3+q^2 &= q^2 (1+q)^2.
\end{align*}
\end{lemma}

\begin{proof}
Direct expansion:
\begin{align*}
(1+q)^6 &= 1 + 6q + 15q^2 + 20q^3 + 15q^4 + 6q^5 + q^6,\\
q(1+q)^2 (q^2 + 4q + 1) &= q(1 + 2q + q^2)(q^2 + 4q + 1)\\
&= q\bigl((q^2+4q+1) + 2q(q^2+4q+1) + q^2(q^2+4q+1)\bigr)\\
&= q\bigl(q^2+4q+1 + 2q^3+8q^2+2q + q^4+4q^3+q^2\bigr)\\
&= q\bigl(q^4 + 6q^3 + 10q^2 + 6q + 1\bigr)\\
&= q^5 + 6q^4 + 10q^3 + 6q^2 + q,\\
q^2 (1+q)^2 &= q^2 (1 + 2q + q^2) = q^2 + 2q^3 + q^4.\qedhere
\end{align*}
\end{proof}

\begin{lemma}[Statistic table]\label{lem:nlambda}
For $\lambda \vdash 4$: $n((4))=0$, $n((3,1))=1$, $n((2,2))=2$,
$n((2,1,1))=3$, $n((1^4))=6$.
\end{lemma}

\begin{proof}
$n(\lambda)=\sum_i (i-1)\lambda_i$; substitute.
\end{proof}

\begin{lemma}[Palindromy and nonnegativity of $p_\lambda$]\label{lem:pallnn}
For each $\lambda \in \{(4),(3,1),(2,2)\}$, the residual polynomial
$p_\lambda(q) := e_\lambda(q)/(q^{n(\lambda)}(1+q)^2)$ is a palindromic
polynomial in $\mathbb{Z}_{\geq 0}[q]$ of degree $d_\lambda$:
\[
p_{(4)}(q) = 1+4q+6q^2+4q^3+q^4, \quad
p_{(3,1)}(q) = 1+4q+q^2, \quad
p_{(2,2)}(q) = 1.
\]
\end{lemma}

\begin{proof}
Division: using Lemma~\ref{lem:factor},
\begin{align*}
\frac{e_{(4)}}{q^0(1+q)^2} &= \frac{(1+q)^6}{(1+q)^2} = (1+q)^4,\\
\frac{e_{(3,1)}}{q^1(1+q)^2} &= \frac{q(1+q)^2(q^2+4q+1)}{q(1+q)^2} = q^2+4q+1,\\
\frac{e_{(2,2)}}{q^2(1+q)^2} &= \frac{q^2(1+q)^2}{q^2(1+q)^2} = 1.
\end{align*}

Palindromy: a polynomial $p(q) = \sum_{k=0}^{d} c_k q^k$ is palindromic if
$c_k = c_{d-k}$ for all $k$. Coefficients:
$(1+q)^4 = 1,4,6,4,1$ (palindromic);
$q^2+4q+1$ has coefficients $1,4,1$ (palindromic);
$1$ is trivially palindromic. All coefficients are nonnegative integers.
\end{proof}

\begin{proof}[Proof of Theorem~\ref{thm:main}]
Lemmas~\ref{lem:charpoly}--\ref{lem:pallnn} establish each entry of the
table:
\begin{itemize}
\item Column $\dim V_\lambda$: standard.
\item Column $n(\lambda)$: Lemma~\ref{lem:nlambda}.
\item Column rank: Lemma~\ref{lem:charpoly}.
\item Column $e_\lambda(q)$: Lemma~\ref{lem:charpoly} + Lemma~\ref{lem:factor}.
\item Column $p_\lambda(q)$: Lemma~\ref{lem:pallnn}.
\end{itemize}
Identity~\eqref{eq:rosetta} for each $\lambda$ with $\rho_\lambda(\Pi_q) \ne 0$
follows from the explicit factorizations in Lemma~\ref{lem:factor}.
Palindromy and nonnegativity of $p_\lambda$ are Lemma~\ref{lem:pallnn}.
\end{proof}

\section{Structural reading}

\subsection{The three factors as three structural sources}

The decomposition $e_\lambda = q^{n(\lambda)} \cdot (1+q)^{\lfloor n/2 \rfloor}
\cdot p_\lambda$ separates $e_\lambda$ into three pieces, each carrying a
distinct structural provenance.

\paragraph{The $q^{n(\lambda)}$ factor.}
$n(\lambda)$ is the cocharge statistic on standard Young tableaux of shape
$\lambda$, equal to the minor index of the superstandard tableau; it appears
as the $q$-Kostka exponent in $\tilde{K}_{\lambda\lambda}(q) = q^{n(\lambda)}$.
The occurrence of $q^{n(\lambda)}$ as a factor of $e_\lambda$ is consistent
with $\Pi_q$ acting on $V_\lambda$ as a weighted sum of Young-seminormal
matrix units where the weight records the lower-triangular word statistic.

\paragraph{The $(1+q)^{\lfloor n/2\rfloor}$ factor.}
For $n=4$, $\lfloor n/2 \rfloor = 2$. This is the \emph{descent-paired piece}:
one factor of $(1+q)$ per descent-pair block $(2i-1, 2i)$. In the rank
hierarchy proof (\texttt{project\_rank\_hierarchy.md}), the image of
$\Pi_q$ on the regular representation has rank exactly $|D_n| =
n!/2^{\lfloor n/2\rfloor}$, and this factor is the one that encodes the
$\lfloor n/2\rfloor$ pair-reflection relations $R_{2i-1} = R_{2i}$.

\paragraph{The palindromic residual $p_\lambda(q)$.}
The residual $p_\lambda(q)$ is the Kazhdan--Lusztig piece. Palindromy of
KL-type polynomials reflects self-duality under the Kazhdan--Lusztig bar
involution $\overline{T_w} = T_{w^{-1}}^{-1}$; Woo's slogan (MO~Q477404,
memory reference \texttt{reference\_woo\_kl\_slogan.md}) is that
\emph{KL polynomials are the correction factor for failure of palindromy}.
The observed palindromy of $p_\lambda$ is the analogue slogan for the
eigenvalue Rosetta Stone: after factoring out the cocharge and the
descent-paired pieces, what remains is intrinsically palindromic.

\subsection{The rank-hierarchy selection rule}

The zero-rank rows $\lambda \in \{(2,1,1), (1^4)\}$ are predicted by the
rank-hierarchy criterion
(\texttt{project\_rank\_hierarchy.md}, \texttt{project\_frobenius\_proof.md}):
$\rho_\lambda(\Pi_q) \ne 0$ iff $\ell(\lambda) \leq (n+1)/2$.

\begin{proposition}[Selection rule at $n=4$]
For $\lambda \vdash 4$, $\rho_\lambda(\Pi_q) \ne 0$ iff $\ell(\lambda) \leq 2$.
\end{proposition}

\begin{proof}
$(n+1)/2 = 2.5$, so the criterion is $\ell(\lambda) \leq 2$. The partitions
of $4$ of length $\leq 2$ are $(4)$, $(3,1)$, $(2,2)$; those of length
$\geq 3$ are $(2,1,1)$, $(1^4)$. Matches the table of Theorem~\ref{thm:main}.
\end{proof}

\subsection{The $6 = |D_4|$ coincidence}\label{sec:six}

The total number of nonzero eigenvalues of $M_{\Pi_q}$ with multiplicity is
\[
\underbrace{1}_{(4)} + \underbrace{3}_{(3,1)} + \underbrace{2}_{(2,2)}
\;=\; 6 \;=\; |D_4|,
\]
since each surviving irrep $V_\lambda$ contributes $\dim V_\lambda$ copies
and $\rho_\lambda(\Pi_q)$ has rank $1$ on each copy. This matches the
rank of $M_{\Pi_q}$ on the regular representation (companion note
\texttt{2026-04-22-bms-direct-hypothesis-refuted.tex},
Proposition on rank of $\Pi_q$) and the descent-paired-set cardinality
$|D_4|=6$.

In Elias--Williamson theory (memory
\texttt{project\_soergel\_convergence.md}), $\Pi_q = BS(\mathbf{i})$
decomposes in the Karoubi completion as
\[
BS(\mathbf{i}) \;=\; B_{w_0} \;\oplus\; \bigoplus_{w < w_0} m_w\, B_w
\]
with $m_w \in \mathbb{Z}_{\geq 0}$ the $\mathbf{i}$-Bott--Samelson
multiplicities. The number of distinct indecomposable summands (counted
with multiplicity) equals the rank of the induced map on the
Grothendieck group of the realization.

\begin{conjecture}[One nonzero eigenvalue per Soergel summand]\label{conj:soergel-eig}
For $n=4$ and the staircase reduced word $\mathbf{i} = s_1 s_2 s_1 s_3 s_2 s_1$,
the number of indecomposable Soergel summands of $BS(\mathbf{i})$ in the
Karoubi completion of $\mathcal{H}_{\mathrm{EW}}(S_4)$, counted with
multiplicity $m_w$, equals $|D_4| = 6$, and coincides with the number of
nonzero eigenvalues of $\Pi_q$ on the regular representation. More
speculatively, each Soergel summand $B_w$ with $m_w > 0$ contributes a
single nonzero eigenvalue to the spectrum of $\Pi_q$ on the appropriate
Specht module, and the residual polynomial $p_\lambda$ records the
$B_w$-multiplicities weighted by character values of $w$ on $V_\lambda$.
\end{conjecture}

\begin{remark}
Conjecture~\ref{conj:soergel-eig} is stated as a conjecture because the
Bott--Samelson multiplicities $m_w$ for the staircase word at $n=4$ are not
computed in this note, and the proposed character-weighted interpretation
of $p_\lambda$ is not verified. The numerical coincidence
$6 = |D_4|$ is observed; the categorical reading is proposed, not proved.
\end{remark}

\section{Computational reproducibility}

All data in Theorem~\ref{thm:main} comes from explicit symbolic computation
in SymPy:

\begin{itemize}
\item \texttt{/home/clio/projects/bms-test/build\_factors.py} constructs
$F_i$ as $24\times 24$ matrices over $\mathbb{Z}[q]$ via the Hecke
right-action on the $T$-basis.

\item \texttt{/home/clio/projects/bms-test/compare\_pi\_bms.py} forms
$M_{\Pi_q} = F_1 F_2 F_1 F_3 F_2 F_1$, verifies rank $6$ at
$q \in \{1,2,3,5\}$, and projects onto the five Specht isotypic components
via central idempotents computed from the Young seminormal form.

\item \texttt{/home/clio/projects/bms-test/followup\_structure.py}
computes the symbolic characteristic polynomial of $M_{\Pi_q}$ and factors
it, producing the five linear factors in Lemma~\ref{lem:charpoly}.

\item Output pickles: \texttt{pi\_bms\_result.pickle} (fields
\texttt{Pi\_correct}, \texttt{per\_irrep\_ranks\_q1},
\texttt{eigs\_q1}) and \texttt{structure\_result.pickle} (field
\texttt{charpoly\_factored}).
\end{itemize}

\begin{thebibliography}{9}

\bibitem{EW} B.~Elias and G.~Williamson.
\emph{Soergel calculus}. Represent. Theory 20 (2016), 295--374.

\bibitem{LLT} A.~Lascoux, B.~Leclerc, and J.-Y.~Thibon.
\emph{The plactic monoid}, in \emph{Algebraic Combinatorics on Words} (M.~Lothaire, ed.),
Cambridge Univ. Press, 2002.

\bibitem{Mac} I.G.~Macdonald.
\emph{Symmetric Functions and Hall Polynomials}, 2nd ed., Oxford Univ. Press, 1995.
(Definition of $n(\lambda)$, eq.~(I.1.5).)

\bibitem{Soe} W.~Soergel.
\emph{The combinatorics of Harish-Chandra bimodules}. J.\ reine angew.\ Math.\
429 (1992), 49--74.

\bibitem{WooMO} A.~Woo.
\emph{MathOverflow Q477404: ``What do Kazhdan--Lusztig polynomials measure?''}
(2024).

\end{thebibliography}

\end{document}
